% DeepGraph-GBM preprint manuscript (bioRxiv style)
\documentclass[11pt,letterpaper]{article}

% ---------- Packages ----------
\usepackage[utf8]{inputenc}
\usepackage[T1]{fontenc}
\usepackage{newtxtext,newtxmath}      % Times-like fonts (bioRxiv look)
\usepackage{courier}                  % Typewriter font fallback for \texttt
\usepackage[top=0.85in,bottom=0.85in,left=0.9in,right=0.9in]{geometry}
\usepackage{graphicx}
\usepackage{float}
\usepackage{caption}
\usepackage{subcaption}
\usepackage{booktabs}
\usepackage{amsmath}
\usepackage{microtype}
\usepackage{xcolor}
\usepackage[numbers,sort&compress,super]{natbib}
\usepackage{authblk}
\usepackage{setspace}
\usepackage{titlesec}
\usepackage{fancyhdr}
\usepackage{lastpage}
\usepackage[hidelinks]{hyperref}
\usepackage{iftex}
\ifXeTeX\renewcommand{\textsuperscript}[1]{\ensuremath{^{\mbox{\fontsize{7}{8}\selectfont #1}}}}\fi

% ---------- bioRxiv-style formatting ----------
\captionsetup{font=small,labelfont={bf},labelsep=period,skip=6pt}
\setlength{\parskip}{0.35em}
\setlength{\parindent}{0pt}
\renewcommand{\baselinestretch}{1.12}

\titleformat{\section}{\large\bfseries}{\thesection}{0.8em}{}
\titleformat{\subsection}{\normalsize\bfseries}{\thesubsection}{0.8em}{}
\titlespacing*{\section}{0pt}{1.0em}{0.4em}
\titlespacing*{\subsection}{0pt}{0.8em}{0.3em}

% header / footer
\pagestyle{fancy}
\fancyhf{}
\fancyfoot[C]{\small \thepage\ of \pageref{LastPage}}
\renewcommand{\headrulewidth}{0pt}

% niche colors (Okabe-Ito) for inline reference
\definecolor{cold}{HTML}{D55E00}
\definecolor{hot}{HTML}{0072B2}
\definecolor{normal}{HTML}{009E73}
\definecolor{necrotic}{HTML}{CC79A7}

\graphicspath{{figures/}}

% ---------- Title & Authors ----------
\title{\vspace{-1.5em}\bfseries DeepGraph-GBM: Spatial Graph Structure Does Not Improve Supervised Niche Prediction in Glioblastoma Visium Data --- a Patient-Level Benchmark}

\author[1,2,*]{Md. Jubayer Hossain}
\affil[1]{Center for Health Innovation, Research, Action, and Learning - Bangladesh}
\affil[2]{DeepBio}
\affil[*]{Corresponding author: \texttt{hossainmj.me@gmail.com}}
\date{}


\renewcommand\Affilfont{\small}
\setlength{\affilsep}{0.5em}

\begin{document}
\maketitle
\thispagestyle{fancy}

% ============================================================
\begin{abstract}
\small
\noindent\textbf{Background.} Glioblastoma (GBM) is organized into recurring spatial niches (necrotic and perinecrotic zones, cellular tumor dominated by mesenchymal-like (MES) cancer cells and immunosuppressive myeloid cells, rarer lymphocyte-infiltrated regions, and the infiltrated brain edge) that differ in biology and are linked to prognosis. Spatial transcriptomics resolves these niches, and the physical adjacency of Visium spots makes graph neural networks (GNNs) a natural modeling choice. Whether neighborhood aggregation actually improves supervised niche prediction over per-spot expression has not been tested under controlled conditions.

\noindent\textbf{Methods.} We built DeepGraph-GBM, an open multi-task framework that represents each Visium section as a spatial $k$-nearest-neighbor graph and trains a shared encoder with three heads: four-class niche classification (immune-cold, immune-hot, normal, necrotic), spot-level MES-state regression, and a slide-level risk score. We trained on 100{,}488 spots (75{,}697 with rule-derived niche labels) from 28 sections of 13 IDH-wildtype GBM patients in the SNUH spatial atlas and compared a GraphSAGE GNN, a four-head graph attention network (GAT), a multilayer perceptron (MLP) and multinomial logistic regression under a matched training and scoring protocol, on held-out patients, across 6 patient-level splits and 3 seeds (72 trained models). External transfer was assessed on 13 Visium samples from an independent cohort, and a niche-signature risk score was tested against overall survival in TCGA-GBM.

\noindent\textbf{Results.} On the primary split, the GraphSAGE GNN reached a test macro-F1 of $0.489 \pm 0.027$ (mean $\pm$ SD over seeds) with one-vs-rest AUROCs of 0.74--0.95. Graph-free models trained and scored identically did as well or better: MLP $0.578 \pm 0.021$ and logistic regression $0.554$. An apparent GNN advantage in our earlier analysis disappeared once all models were scored with the same validation-tuned thresholds. Across six patient-level splits, logistic regression exceeded GraphSAGE on every split (mean difference $0.086$, paired $p = 4 \times 10^{-4}$). GAT improved on GraphSAGE ($+0.037$, $p = 0.007$) but did not exceed the MLP ($p = 0.89$). The choice of held-out patients dominated performance: across-split variance was 3.5--8.1$\times$ the seed variance, and split-mean macro-F1 ranged from 0.32 to 0.70. Most errors came from normal and necrotic spots being labeled immune-cold. In the external cohort, predictions were almost entirely necrotic or immune-hot (immune-cold $< 1\%$ in 11 of 13 samples), showing strong cohort shift. Predicted MES probability correlated with an independent MES signature in bulk-tumor samples (mean $r = 0.61$) but not in infiltrating-edge samples (mean $r = 0.04$). The niche-signature risk score did not stratify TCGA-GBM survival ($n = 166$; HR 1.12, 95\% CI 0.92--1.37; C-index 0.52).

\noindent\textbf{Conclusions.} In this setting, spatial graph structure did not improve supervised niche prediction over per-spot expression, and the identity of the held-out patients mattered more than the choice of architecture. We release the full pipeline, splits and results as a benchmark, and recommend matched scoring, baselines without graph structure, and repeated patient-level splits as minimum standards for spatial deep-learning claims.

\vspace{0.5em}
\noindent\textbf{Keywords:} glioblastoma; spatial transcriptomics; Visium; graph neural network; benchmarking; tumor microenvironment; mesenchymal state; reproducibility
\end{abstract}

\clearpage

% ============================================================
\section{Introduction}

Glioblastoma, IDH-wildtype, is the most common malignant primary brain tumor in adults \cite{ostrom2023,louis2021}. Median survival remains around 15 months despite maximal resection, radiotherapy and temozolomide \cite{stupp2005}. Bulk expression profiling first separated GBM into proneural, classical and mesenchymal subtypes \cite{verhaak2010,tcga2008}. Single-cell studies then showed that individual tumors contain malignant cells in several plastic states (neural-progenitor-like, oligodendrocyte-progenitor-like, astrocyte-like and mesenchymal-like) whose proportions vary within and between patients \cite{neftel2019}. The mesenchymal (MES) program is of particular interest because it is associated with tumor-associated macrophages and microglia \cite{wang2017}. Macrophage-derived signals can drive cancer cells into MES-like states \cite{hara2021}. Myeloid cells make up a large share of the immune compartment in GBM, while effector T cells are scarce \cite{klemm2020,pomboantunes2021,khan2023,decordova2020,lin2024}.

These cell states are spatially organized. The Ivy Glioblastoma Atlas Project defined transcriptionally distinct anatomic features, including the leading edge, infiltrating tumor, cellular tumor, pseudopalisading cells around necrosis and microvascular proliferation \cite{puchalski2018}. Spatial transcriptomics \cite{stahl2016,moses2022} has since shown that GBM is arranged into recurring niches shaped by hypoxia, tumor--host interactions and immune composition \cite{ravi2022,greenwald2024,ruizmoreno2025}. Greenwald et al.\ described a multi-layered organization in which hypoxia-associated MES programs structure the surrounding tissue \cite{greenwald2024}. The recent SNUH atlas integrated Visium with matched single-cell and protein data to characterize the cellular composition of each anatomical niche \cite{park2026}. Together, these studies suggest that a tissue section could be summarized as a map of niches that differ in immune state and prognosis.

Producing such maps currently requires expert annotation, deconvolution against single-cell references \cite{kleshchevnikov2022,cable2022,li2022bench}, or unsupervised spatial clustering followed by manual interpretation. Many spatial clustering methods treat a Visium section as a graph of adjacent spots, including SpaGCN \cite{hu2021spagcn}, STAGATE \cite{dong2022stagate}, GraphST \cite{long2023graphst} and related approaches \cite{zhao2021bayesspace,xu2023spacel,peng2023,liu2023review,palla2022}. The underlying assumption is that a spot's identity depends on its neighbors, so aggregating neighborhood information with GraphSAGE \cite{hamilton2017} or graph attention \cite{velickovic2018} should improve on per-spot expression. Benchmarks of unsupervised spatial clustering show that method rankings depend strongly on the dataset \cite{yuan2024bench}. In graph machine learning more broadly, simple or graph-free models often match more complex GNNs once training procedures and data splits are controlled \cite{shchur2018,errica2020}. Mean aggregation can also over-smooth node representations \cite{li2018deeper} and performs poorly when adjacent nodes carry different labels \cite{zhu2020homophily}. Both conditions plausibly arise at niche boundaries in GBM.

Whether graph structure helps \emph{supervised} prediction of GBM niches has therefore not been tested directly, and the answer depends on evaluation choices. With few patients, results can depend on which patients are held out \cite{bouthillier2021,varoquaux2022}. Applying different decision rules to different models, or tuning on information that is not available at test time, can also change apparent rankings \cite{kapoor2023}.

Here we present DeepGraph-GBM, an open multi-task framework for spot-level niche classification, MES-state regression and slide-level risk scoring from Visium data, together with a controlled benchmark. Our contributions are:
\begin{enumerate}
  \item \textbf{A matched comparison of graph-based and graph-free models.} GraphSAGE, GAT, an MLP and logistic regression were trained with the same features, optimizer, class weighting and early stopping, and scored with the same validation-tuned thresholds on held-out patients.
  \item \textbf{An assessment of patient-sampling variance.} We repeated the benchmark across six patient-level splits and three seeds (72 models) and partitioned performance variance into split and seed components.
  \item \textbf{A demonstration that scoring rules can reverse model rankings.} The GNN appeared superior in an earlier analysis only because it was scored differently from the baselines.
  \item \textbf{External and clinical validation, including negative results.} We report cross-cohort predictions on an independent Visium cohort and a survival analysis in TCGA-GBM. The survival association was not significant.
\end{enumerate}

% ============================================================
\section{Results}

\subsection{Cohort, labels and model}

We used the IDH-wildtype GBM sections of the SNUH spatial atlas \cite{park2026}: 28 Visium sections (23 FFPE, 5 fresh-frozen) from 13 patients, comprising 100{,}488 spots (Supplementary Table~\ref{tab:s3}). Niche labels were derived from the atlas's per-spot anatomical-feature annotations (Ivy GAP-style categories \cite{puchalski2018}) and metaprogram annotations using priority-ordered rules (Methods). We refer to these as \emph{silver-standard} labels. Spots were labeled \emph{necrotic} (pseudopalisading or perinecrotic features), \emph{immune-hot} (immune metaprogram outside necrotic features), \emph{normal} (concordant leading-edge grey or white matter) or \emph{immune-cold} (cellular or infiltrating tumor features with tumor, infiltrating or hypoxic metaprograms). Vascular and discordant spots were excluded from classification. In total, 75{,}697 spots were labeled. The class distribution varied strongly by patient. In the primary split, the three test patients (SNU16, SNU33, SNU34) contributed 16{,}377 labeled spots: 11{,}125 immune-cold (67.9\%), 185 immune-hot (1.1\%), 1{,}455 normal and 3{,}612 necrotic. The two validation patients contributed only 15 immune-hot spots.

Each section was represented as a $k$-nearest-neighbor graph ($k = 6$) over spot coordinates, with 3{,}000 highly variable genes as node features. A three-layer encoder (3000$\rightarrow$256$\rightarrow$256$\rightarrow$128) fed three heads: a four-class niche classifier, a sigmoid MES-probability regressor trained on the fraction of MES malignant states per spot, and an attention-pooled slide-level risk head trained on a niche-composition pseudo-risk (Methods). In the matched benchmark, the MLP received the same features without edges, so the comparison isolates the effect of neighborhood aggregation.

\subsection{Learned embeddings and per-class discrimination}

Figure~\ref{fig:umap} shows a UMAP \cite{mcinnes2018} projection of GraphSAGE embeddings for all spots in the held-out test sections, including spots excluded from labeling. Immune-cold spots form several large, separate clusters, and necrotic spots form a dense region that partly overlaps a normal-enriched cluster. Immune-hot and normal spots lie mostly in a central area occupied largely by excluded spots. The presence of several separate immune-cold clusters suggests that the embedding also reflects section- or patient-level differences, in addition to niche identity.

\begin{figure}[H]
  \centering
  \includegraphics[width=0.55\textwidth]{Figure1_umap_niche.png}
  \caption{\textbf{GraphSAGE spot embeddings for held-out test sections.} UMAP projection of 128-dimensional encoder embeddings for all spots from test patients SNU16, SNU33 and SNU34 (representative seed-7 model), colored by silver-standard niche label. Grey points are spots excluded from labeling (vascular or discordant annotations).}
  \label{fig:umap}
\end{figure}

With per-class thresholds tuned on validation patients, the representative seed-7 GraphSAGE model achieved a test macro-F1 of 0.504, weighted-F1 of 0.698 and accuracy of 0.740 (Supplementary Table~\ref{tab:s8}). One-vs-rest AUROC was 0.82 for immune-cold, 0.98 for immune-hot, 0.73 for normal and 0.96 for necrotic (Figure~\ref{fig:roc}). Because immune-hot accounted for only 1.1\% of test spots, AUROC overstates how well this class is resolved. The corresponding area under the precision--recall curve (AUPRC) was 0.57, against a chance level of 0.011 \cite{saito2015}. For normal spots, AUPRC was 0.27 against a chance level of 0.089. Test macro-F1 was lower than validation macro-F1 (0.504 vs.\ 0.584), consistent with the different class composition of the validation and test patients.

\begin{figure}[H]
  \centering
  \includegraphics[width=0.52\textwidth]{Figure2_roc_curves.png}
  \caption{\textbf{Per-class discrimination on held-out patients.} One-vs-rest ROC curves for the four niche classes (representative seed-7 GraphSAGE model, test patients SNU16, SNU33 and SNU34; $n = 16{,}377$ labeled spots). AUROC is shown in the legend. Because immune-hot accounts for 1.1\% of test spots, see Supplementary Table~\ref{tab:s8} for precision--recall metrics.}
  \label{fig:roc}
\end{figure}

The confusion matrix shows the main error pattern (Figure~\ref{fig:confusion}). Under default argmax decisions, 92\% of immune-cold spots were classified correctly, but most normal (80\%) and necrotic (82\%) spots were also assigned to immune-cold, and 83\% of immune-hot spots were assigned to normal. Validation-tuned thresholds changed this substantially. With thresholds of 0.30 (immune-cold), 0.05 (immune-hot), 0.725 (normal) and 0.15 (necrotic), test recall was 0.95 for immune-cold, 0.94 for immune-hot, 0.31 for necrotic and 0.15 for normal. Normal-to-immune-cold and necrotic-to-immune-cold confusion was therefore the main remaining error. Because the immune-hot threshold was tuned on only 15 validation spots, the model's immune-hot performance depends heavily on that tuning.

\begin{figure}[H]
  \centering
  \includegraphics[width=0.52\textwidth]{Figure3_confusion_matrix.png}
  \caption{\textbf{Confusion matrix under argmax decisions.} Row-normalized confusion matrix for the representative seed-7 GraphSAGE model on held-out test patients, using default argmax over softmax probabilities. Rows are silver-standard labels; columns are predictions. Quantitative metrics in the text and tables use validation-tuned per-class thresholds, which shift errors substantially (see text).}
  \label{fig:confusion}
\end{figure}

\subsection{Graph-free baselines match or exceed the GNN under matched evaluation}

To test whether neighborhood aggregation adds predictive signal, we compared GraphSAGE with a four-head GAT encoder, an MLP and multinomial logistic regression. All models used identical inputs and splits, the same optimizer, cosine learning-rate schedule, boosted class weights and early stopping on validation macro-F1, and the same validation-tuned per-class thresholds applied to test (Methods). On the primary split over three seeds (Table~\ref{tab:comparison}; Supplementary Table~\ref{tab:s6}), the MLP achieved the highest macro-F1 ($0.578 \pm 0.021$), followed by logistic regression ($0.554$), GAT ($0.539 \pm 0.063$) and GraphSAGE ($0.489 \pm 0.027$). The MLP exceeded GraphSAGE on every seed (mean paired difference $-0.089$, 95\% CI $-0.105$ to $-0.073$; paired $t = -24.1$, df $= 2$, $p = 0.0017$). The comparison with logistic regression did not reach significance ($p = 0.053$). The MLP also had the highest MES AUROC ($0.826$ vs.\ $0.751$ for GraphSAGE) and MES Pearson correlation ($0.459$ vs.\ $0.359$). GraphSAGE was competitive only for necrotic AUROC (0.951 vs.\ 0.943 for the MLP).

These seed-level tests are limited. With three seeds the Wilcoxon test cannot reach $p < 0.25$, and all seeds share one patient split, so the $p$-values reflect training-seed variability rather than patient sampling. An independent re-run of the matched MLP on the same split within the multi-split pipeline gave $0.571 \pm 0.021$, indicating run-to-run differences of about 0.007 macro-F1. We therefore treat the multi-split analysis below as the primary test.

\begin{table}[H]
\centering
\caption{\textbf{Matched comparison on the primary patient split.} Mean $\pm$ SD over seeds 42, 7 and 123. All models were trained under the same optimization protocol and scored with per-class thresholds tuned on validation patients (SNU24, SNU51) and applied to test patients (SNU16, SNU33, SNU34). Logistic regression is deterministic, so its SD is zero. Paired $t$-tests on per-seed macro-F1 (df $= 2$): GraphSAGE vs.\ MLP $p = 0.0017$; vs.\ logistic $p = 0.053$; vs.\ GAT $p = 0.41$. Best value per row in bold.}
\label{tab:comparison}
\small
\begin{tabular}{lcccc}
\toprule
\textbf{Metric} & \textbf{GraphSAGE} & \textbf{GAT} & \textbf{MLP} & \textbf{Logistic} \\
\midrule
Macro-F1            & 0.489 $\pm$ 0.027 & 0.539 $\pm$ 0.063 & \textbf{0.578 $\pm$ 0.021} & 0.554 \\
Weighted-F1         & 0.707 $\pm$ 0.071 & 0.738 $\pm$ 0.051 & \textbf{0.764 $\pm$ 0.008} & 0.724 \\
Accuracy            & 0.750 $\pm$ 0.050 & 0.742 $\pm$ 0.064 & \textbf{0.789 $\pm$ 0.004} & 0.748 \\
AUROC immune-cold   & 0.782 $\pm$ 0.108 & 0.805 $\pm$ 0.069 & \textbf{0.838 $\pm$ 0.012} & 0.760 \\
AUROC immune-hot    & 0.947 $\pm$ 0.051 & 0.953 $\pm$ 0.018 & 0.987 $\pm$ 0.002 & \textbf{0.991} \\
AUROC normal        & 0.739 $\pm$ 0.051 & 0.760 $\pm$ 0.074 & \textbf{0.829 $\pm$ 0.017} & 0.687 \\
AUROC necrotic      & \textbf{0.951 $\pm$ 0.007} & 0.938 $\pm$ 0.008 & 0.943 $\pm$ 0.007 & 0.925 \\
MES AUROC           & 0.751 $\pm$ 0.041 & 0.668 $\pm$ 0.076 & \textbf{0.826 $\pm$ 0.016} & 0.751 \\
MES Pearson $r$     & 0.359 $\pm$ 0.016 & 0.273 $\pm$ 0.068 & \textbf{0.459 $\pm$ 0.021} & 0.410 \\
\bottomrule
\end{tabular}
\end{table}

\subsection{Mismatched scoring reverses the model ranking}

In an earlier version of this analysis, the GNN was scored with validation-tuned thresholds, whereas the baselines were trained for a fixed 120 epochs without early stopping and scored with argmax. Under that protocol, the GNN appeared best (macro-F1 0.489 vs.\ 0.449 for the MLP and 0.350 for logistic regression; Figure~\ref{fig:reversal}A; Supplementary Table~\ref{tab:s1}). Under the matched protocol, the ranking reversed (Figure~\ref{fig:reversal}B).

Two controls separate the contributions of scoring and training (Figure~\ref{fig:reversal}C). Logistic regression is trained identically under both protocols, so rescoring it isolates the effect of the decision rule. Switching from argmax to tuned thresholds raised its macro-F1 from 0.350 to 0.554, a gain of 0.20. For the matched-protocol MLP, the change in decision rule had a smaller effect (argmax $0.557 \pm 0.050$ vs.\ tuned $0.571 \pm 0.021$). Most of the original MLP deficit was therefore due to the 120-epoch training schedule without early stopping, not to the decision rule. For both baselines, the apparent GNN advantage came from evaluation differences rather than from graph structure.

\begin{figure}[H]
  \centering
  \includegraphics[width=\textwidth]{Figure_protocol_reversal.png}
  \caption{\textbf{Matched evaluation reverses the apparent GNN advantage.} (A) Original, mismatched evaluation: GraphSAGE scored with validation-tuned thresholds (mean $\pm$ SD, three seeds); baselines trained for 120 epochs without early stopping and scored by argmax (single run). (B) Matched protocol: all models trained identically and scored with tuned thresholds (mean $\pm$ SD, three seeds). (C) Effect of the decision rule alone: deterministic logistic regression and the matched-protocol MLP (mean of three seeds) scored by argmax and by tuned thresholds.}
  \label{fig:reversal}
\end{figure}

\subsection{Patient-split identity dominates benchmark variance}

To estimate patient-sampling variability, we repeated the full matched benchmark on six patient-level splits (8 training, 2 validation and 3 test patients; six distinct test sets, split 0 being the primary split) with three seeds each, for all four models (72 trained models; Figure~\ref{fig:robustness}; Supplementary Tables~\ref{tab:s7}, \ref{tab:s9}).

Split identity was the largest source of variance. Split-mean macro-F1 ranged from 0.32 to 0.63 for GraphSAGE, 0.36 to 0.70 for GAT, 0.37 to 0.69 for the MLP and 0.43 to 0.68 for logistic regression, and all four models rose and fell together across splits (Figure~\ref{fig:robustness}A). For the stochastic models, the variance of split means was 3.5$\times$ (GAT), 3.6$\times$ (GraphSAGE) and 8.1$\times$ (MLP) the mean within-split seed variance (SD across splits 0.11--0.12 vs.\ mean within-split SD 0.04--0.06). Logistic regression has no seed variance, so this ratio is undefined for it. The gap between the best and worst split (about 0.3 macro-F1) was roughly three times the largest difference between models.

Averaged across splits, logistic regression had the highest macro-F1 ($0.572 \pm 0.090$), followed by the MLP ($0.526 \pm 0.108$), GAT ($0.524 \pm 0.116$) and GraphSAGE ($0.487 \pm 0.109$). Logistic regression exceeded GraphSAGE on all six splits (mean paired difference $0.086$, 95\% CI $0.059$--$0.112$; $t = 8.27$, df $= 5$, $p = 4.2 \times 10^{-4}$; Wilcoxon $p = 0.031$, the minimum attainable with six pairs). The MLP exceeded GraphSAGE on five of six splits, but the difference was not significant ($0.040$, 95\% CI $-0.018$ to $0.098$; $p = 0.14$). The fixed-split result above therefore overstated the certainty of the MLP advantage. GAT exceeded GraphSAGE on all six splits ($+0.037$, $p = 0.007$), so learned attention over neighbors performed better than mean aggregation. However, GAT was indistinguishable from the MLP ($-0.003$, $p = 0.89$) and below logistic regression ($-0.049$, $p = 0.034$). Thus the better graph encoder narrowed the gap but did not outperform graph-free models.

\begin{figure}[H]
  \centering
  \includegraphics[width=\textwidth]{Figure_split_robustness.png}
  \caption{\textbf{Benchmark variance across six patient-level splits.} (A) Test macro-F1 for GraphSAGE, MLP, logistic regression and GAT on splits S0--S5 (S0 is the primary split). Large markers and lines show split means over three seeds; small markers show individual seeds. (B) Paired macro-F1 difference of each model relative to GraphSAGE per split (mean $\pm$ SD over seeds). Logistic regression and GAT exceed GraphSAGE on all six splits; the MLP does so on five of six.}
  \label{fig:robustness}
\end{figure}

\subsection{Spatial maps show over-assignment and smoothing}

Figure~\ref{fig:spatial} and Supplementary Figures~\ref{fig:s1}--\ref{fig:s4} compare silver-standard and predicted maps for each test section (argmax decisions, seed-7 model). The maps illustrate the main failure modes more clearly than the summary metrics.

In SNU16A (Figure~\ref{fig:spatial}a--b), the silver-standard map contains a large contiguous necrotic region in the lower-left and many normal spots in the upper half. The predicted map assigns nearly the whole section to immune-cold, with small normal and necrotic patches, and does not recover the necrotic region. In SNU34A (Supplementary Figure~\ref{fig:s4}), the model correctly predicts that most of the section is immune-cold. It also predicts two contiguous normal regions along the upper edge where the silver standard has only scattered normal spots, and it misses scattered necrotic spots on the left. In SNU33A (Supplementary Figure~\ref{fig:s2}), many spots are unlabeled. Predictions there are mostly normal with contiguous immune-hot patches, whereas silver-standard immune-hot spots are scattered.

Across sections, the predicted maps are spatially smoother and more contiguous than the silver-standard maps. This is consistent with neighborhood averaging \cite{li2018deeper}, which can enforce local consistency but can also remove small or scattered niches. The predicted MES-probability map for SNU16A (Figure~\ref{fig:spatial}c) has its highest values in the lower-left, overlapping the silver-standard necrotic region. This pattern agrees with reports linking MES and hypoxia-associated programs to perinecrotic regions \cite{greenwald2024,neftel2019}, even though the classifier did not recover the necrotic region itself.

\begin{figure}[H]
  \centering
  \begin{subfigure}[b]{0.32\textwidth}
    \includegraphics[width=\textwidth]{Figure4a_SNU16A_niche_true.png}
    \caption{Silver-standard niche}
  \end{subfigure}\hfill
  \begin{subfigure}[b]{0.32\textwidth}
    \includegraphics[width=\textwidth]{Figure4b_SNU16A_niche_pred.png}
    \caption{Predicted niche}
  \end{subfigure}\hfill
  \begin{subfigure}[b]{0.32\textwidth}
    \includegraphics[width=\textwidth]{Figure4c_SNU16A_mes.png}
    \caption{Predicted MES probability}
  \end{subfigure}
  \caption{\textbf{Silver-standard and predicted maps for held-out section SNU16A.} (a) Rule-derived niche labels (grey: excluded). (b) Predicted niches (seed-7 GraphSAGE, argmax), which assign most of the section to immune-cold and miss the lower-left necrotic region. (c) Predicted per-spot MES probability, highest in the lower-left region that overlaps the silver-standard necrotic zone.}
  \label{fig:spatial}
\end{figure}

\subsection{External cohort: strong shift in niche composition, region-dependent MES agreement}

We applied the seed-7 model without retraining to 13 GBM Visium samples (25{,}488 spots) from Greenwald et al.\ \cite{greenwald2024}, annotated by sampling region (bulk, contrast-enhancing, infiltrating, necrotic; Supplementary Table~\ref{tab:s2}). Predicted composition differed sharply from the training cohort. Across all external spots, 56.1\% were predicted necrotic, 40.3\% immune-hot, 2.2\% immune-cold and 1.4\% normal, and 11 of 13 samples had fewer than 1\% immune-cold spots. Immune-cold was the majority class in the SNUH test patients (67.9\%) and is expected to be abundant in bulk tumor \cite{greenwald2024}. Its near-absence therefore indicates cohort shift and miscalibration rather than true biological differences. Predicted composition also did not follow the annotated sampling region. The single necrotic-region sample had a lower predicted necrotic fraction (0.43) than the bulk samples (mean 0.61), and the predicted necrotic fraction in infiltrating samples ranged from 0.08 to 0.79.

Agreement for the MES head depended on region. Per-sample Pearson correlation between predicted MES probability and an independent Neftel MES meta-module score \cite{neftel2019} averaged $0.30$ (median 0.32, range $-0.34$ to $0.65$). It was high in bulk-tumor samples (mean $r = 0.61$, $n = 4$), intermediate in contrast-enhancing samples (0.24, $n = 3$) and near zero in infiltrating-edge samples (0.04, $n = 4$). The pooled spot-level correlation was higher ($r = 0.68$), but it is driven largely by differences between samples and should not be read as within-sample agreement. Thus the MES head partly transferred where malignant cells are dense, whereas the niche classifier did not transfer.

\subsection{Niche-signature risk does not stratify TCGA-GBM survival}

No Visium GBM cohort with matched survival was available. The slide-level head was therefore trained on a niche-composition pseudo-risk, and we evaluated the same biological hypothesis in bulk data. Each TCGA-GBM tumor was scored with fixed niche gene signatures (MES/myeloid, T-cell, hypoxia and normal-brain genes), and these scores were combined with the same weights that defined the pseudo-risk (Methods). This score is derived from the niche definitions, not from the trained network. In 166 expression-matched tumor samples from 165 patients (132 deaths), the score did not stratify overall survival (Figure~\ref{fig:km}; Table~\ref{tab:survival}). The univariate Cox hazard ratio per unit score was 1.12 (95\% CI 0.92--1.37, $p = 0.25$), the concordance index was 0.52 \cite{harrell1996}, the age-adjusted hazard ratio was 1.03 ($p = 0.81$), and the median-split log-rank $p$ was 0.42 (median overall survival 424 vs.\ 444 days). This hypothesis, that bulk niche composition is prognostic, was not supported in this cohort.

\begin{figure}[H]
  \centering
  \includegraphics[width=0.52\textwidth]{Figure5_km_survival.png}
  \caption{\textbf{Overall survival by niche-signature risk in TCGA-GBM.} Kaplan--Meier curves (95\% confidence bands) for tumors above and below the median risk score ($n = 83$ each). Log-rank $p = 0.42$.}
  \label{fig:km}
\end{figure}

\begin{table}[H]
\centering
\caption{\textbf{TCGA-GBM survival analysis of the niche-signature risk score} (166 tumor samples from 165 patients; 132 events).}
\label{tab:survival}
\small
\begin{tabular}{lc}
\toprule
\textbf{Analysis} & \textbf{Result} \\
\midrule
Univariate Cox HR (95\% CI)        & 1.12 (0.92--1.37), $p = 0.25$ \\
Age-adjusted Cox HR                & 1.03, $p = 0.81$ \\
Concordance index                  & 0.52 \\
Kaplan--Meier log-rank $p$ (median split) & 0.42 \\
Median OS, high vs.\ low risk      & 424 vs.\ 444 days \\
\bottomrule
\end{tabular}
\end{table}

% ============================================================
\section{Discussion}

We set out to test whether explicitly modeling spatial adjacency improves supervised niche prediction in GBM Visium data. Under matched training and scoring, it did not. A graph-free MLP and logistic regression matched or exceeded a GraphSAGE GNN, and only an attention-based encoder closed most of the gap. The identity of the held-out patients changed performance more than the choice of model.

\textbf{Why did the graph not help?} We see three non-exclusive explanations. First, at Visium resolution each 55-\textmu m spot already pools several cells, so a spot's own transcriptome carries much of the niche information. The niche axes we predict are largely defined by cell-intrinsic programs, such as hypoxia, MES and myeloid, and T-cell genes \cite{neftel2019,greenwald2024}. Second, mean aggregation smooths representations \cite{li2018deeper} and is poorly suited to settings where adjacent spots have different labels \cite{zhu2020homophily}. Immune-hot spots in this cohort are rare and scattered, and niche boundaries are irregular. The contiguous, over-smoothed predicted maps (Figure~\ref{fig:spatial}) are consistent with this explanation, as is the gain from GAT, which can down-weight dissimilar neighbors. Third, the silver-standard labels are themselves derived from per-spot annotations, which may favour per-spot models. Our results do not show that spatial context is biologically unimportant in GBM. They show that, for this label set, data size and resolution, a standard GNN did not extract additional predictive signal from neighboring spots.

\textbf{Evaluation choices can produce apparent gains.} The ranking reversal shows how easily an architectural advantage can be produced by evaluation choices alone. Tuning thresholds for one model while scoring baselines by argmax, or training baselines without the early stopping given to the proposed model, created an apparent GNN advantage of 0.04--0.14 macro-F1. This is the same order of magnitude as many reported improvements. Similar effects have been documented in graph learning \cite{shchur2018,errica2020} and in machine-learning-based science more broadly \cite{kapoor2023,varoquaux2022}. With 13 patients, the choice of held-out patients produced a spread of about 0.3 macro-F1, larger than any difference between models. A comparison on a single split, even with several seeds, would have reported a significant MLP advantage ($p = 0.0017$) that did not hold across splits ($p = 0.14$). This agrees with Bouthillier et al., who found that data sampling is a major and often ignored source of variance in machine-learning benchmarks \cite{bouthillier2021}. We suggest that spatial-omics method papers report, at minimum, (i) baselines without graph structure trained under the same protocol, (ii) identical decision rules for all models, and (iii) results across several patient-level splits.

\textbf{Cohort shift.} In the external cohort, immune-cold was predicted almost nowhere, and predictions did not follow annotated regions. Cohort differences in tissue processing (the training set was mostly FFPE), sequencing depth, region sampling and normalization probably all contribute. In addition, the thresholds and class priors were learned on a cohort in which immune-cold dominates. The MES head transferred better, but only where malignant cells are dense. These results argue for multi-cohort training, explicit domain adaptation, and caution in interpreting niche fractions predicted on new cohorts.

\textbf{Survival.} The niche-signature risk score did not stratify TCGA-GBM survival. Several factors limit this test: bulk expression averages over niches, the signatures are curated rather than learned, and the cohort provides moderate power. Even so, the result suggests that a simple bulk proxy for niche composition is not a strong prognostic marker. Testing spatial niche composition as a prognostic factor will require spatial cohorts with matched outcomes.

\subsection{Limitations}

(i) \textbf{Silver-standard labels.} Labels are derived from the atlas's own annotations by rules, not from independent pathology review, and 25\% of spots were excluded as vascular or discordant. (ii) \textbf{Extreme class imbalance and validation scarcity.} Immune-hot made up 1.1\% of test spots, and its threshold was tuned on 15 validation spots, so immune-hot metrics are fragile. (iii) \textbf{Small patient cohort.} Thirteen patients limit the number of independent splits. The six splits share training patients, and the paired tests ($n = 6$) have limited power. (iv) \textbf{Qualitative figures use argmax decisions.} The confusion matrix, spatial maps and external predictions use argmax, whereas the quantitative metrics use tuned thresholds, so the figures show a more pessimistic decision rule. (v) \textbf{Limited architecture search.} We tested two message-passing encoders with fixed depth and $k = 6$. Edge weights were computed but not used by either encoder. Graph designs suited to heterophily \cite{zhu2020homophily}, multi-scale neighborhoods, histology integration \cite{hu2021spagcn} or deconvolution-derived features \cite{kleshchevnikov2022} may perform differently. (vi) \textbf{Survival.} The survival analysis tested a signature-based proxy, not the trained network, on bulk data.

\subsection{Conclusions}

For supervised niche prediction in GBM Visium data, per-spot expression with a simple classifier was at least as informative as GraphSAGE or GAT encoders, and the choice of held-out patients mattered more than the architecture. We release DeepGraph-GBM, including code, splits, trained weights and all per-run results, as a reproducible benchmark for future spatial models.

% ============================================================
\section{Methods}

\subsection{Data}

\textbf{Training cohort.} We used the IDH-wildtype GBM Visium sections of the SNUH spatial atlas \cite{park2026}, obtained from the UCSC Cell Browser: 28 sections (23 FFPE, 5 fresh-frozen) from 13 patients, comprising 100{,}488 spots with per-spot anatomical-feature, metaprogram and malignant-state annotations (Supplementary Table~\ref{tab:s3}).

\textbf{External cohort.} We used 13 GBM Visium samples from Greenwald et al.\ \cite{greenwald2024} (Zenodo record 12624108; CC-BY), annotated by sampling region.

\textbf{Survival cohort.} We obtained TCGA-GBM bulk RNA-seq and clinical data from the NCI Genomic Data Commons \cite{tcga2008}. Overall survival combined \texttt{days\_to\_death} for deaths with \texttt{days\_to\_last\_follow\_up} for censored cases. The expression-matched set comprised 166 tumor samples from 165 patients (one patient contributed two aliquots), with 132 deaths.

\subsection{Preprocessing, graphs and labels}

Spots with fewer than 200 detected genes were removed. We selected 3{,}000 highly variable genes with the \texttt{seurat\_v3} method \cite{stuart2019} in scanpy \cite{wolf2018}, applied to counts reconstructed from the log-normalized matrix by \texttt{expm1}. Dispersion-based selection on the normalized matrix mostly returned housekeeping genes and reduced cross-cohort transfer. Each section was converted into a $k$-nearest-neighbor graph ($k = 6$) on spot coordinates. Gaussian edge weights $\exp(-d^2 / 2\tilde{d}^2)$, where $\tilde{d}$ is the median neighbor distance, were computed and stored but were not used by either encoder.

Niche labels were assigned by the first matching rule in the order necrotic $\rightarrow$ immune-hot $\rightarrow$ normal $\rightarrow$ immune-cold. The rules were: necrotic, anatomical feature PAN or PNZ; immune-hot, immune metaprogram and anatomical feature not PAN/PNZ; normal, anatomical feature LE\_GM/LE\_WM with concordant metaprogram LE\_GM/LE\_WM; immune-cold, anatomical feature CT, IT or IGN with metaprogram CT, IT or hypoxic. Spots matching no rule (including vascular features) were labeled \emph{exclude} and masked from the classification loss. The MES target was the fraction of malignant MES states (MES, MES.Hyp, MES.Ast) assigned to each spot.

\subsection{Models}

\textbf{GraphSAGE.} Three SAGEConv layers \cite{hamilton2017} with mean aggregation (3000$\rightarrow$256$\rightarrow$256$\rightarrow$128), each followed by batch normalization, ReLU and dropout ($p = 0.3$), implemented in PyTorch Geometric \cite{fey2019} (1{,}751{,}815 parameters including heads).

\textbf{GAT.} Three GATConv layers \cite{velickovic2018} with four heads and the same dimensions. Hidden layers concatenate four heads of width $d/4$, and the final layer averages the heads. Normalization, dropout and heads were identical to GraphSAGE (985{,}863 parameters).

\textbf{Heads.} (1) A niche classifier (linear layers to four logits). (2) A MES regressor (sigmoid output). (3) A slide-level risk head that applies attention pooling over spot embeddings followed by an MLP producing a scalar.

\textbf{Baselines.} The MLP used the same feature input without edges and the same niche and MES heads. Multinomial logistic regression (scikit-learn, balanced class weights, up to 500 iterations) was fitted on the same features, with a separate logistic model for MES status. Its fit is deterministic, so its results do not vary by seed.

\subsection{Training}

The multi-task loss was $\mathcal{L} = 1.0\,\mathcal{L}_\mathrm{CE}^\mathrm{niche} + 0.5\,\mathcal{L}_\mathrm{BCE}^\mathrm{MES} + 0.5\,\mathcal{L}_\mathrm{MSE}^\mathrm{risk}$. The cross-entropy used balanced class weights, multiplied by an additional 3.0 for immune-hot and 1.5 for necrotic. The pseudo-risk target for each slide was the niche-fraction-weighted sum $1.0\,f_\mathrm{cold} + 0.8\,f_\mathrm{necrotic} - 0.6\,f_\mathrm{hot} - 0.8\,f_\mathrm{normal}$ over labeled spots. We used AdamW (learning rate $10^{-3}$, weight decay $10^{-4}$) with cosine annealing, full-graph batches (one section per step), up to 200 epochs, and early stopping on validation macro-F1 (patience 40) with best-checkpoint restoration. The MLP was trained with the same protocol.

\subsection{Evaluation}

\textbf{Splits.} The primary split (seed 42) used training patients SNU17, 18, 21, 23, 25, 27, 43 and 46, validation patients SNU24 and 51, and test patients SNU16, 33 and 34. Five additional splits (generator seeds 1001--1005) used the same 8/2/3 design, with distinct test sets across all six splits.

\textbf{Decision rule.} For every model, a threshold $t_c$ was tuned for each class $c$ on the validation patients to maximize one-vs-rest F1 (grid of 37 values from 0.05 to 0.95). Test spots were then assigned to $\arg\max_c\, p_c / t_c$ using the validation thresholds unchanged, so test labels were never used for tuning. In the scoring-rule control, the same models were also scored by argmax.

\textbf{Metrics.} We report macro- and weighted-F1, accuracy, one-vs-rest AUROC and AUPRC per class, and, for the MES head, AUROC, AUPRC and Pearson correlation between predicted MES probability and the spot-level MES target.

\textbf{Statistics.} Models were compared with paired $t$-tests and Wilcoxon signed-rank tests on per-seed macro-F1 for the primary split (df $= 2$) and on split-mean macro-F1 across the six splits (df $= 5$). Variance was partitioned into the variance of split means and the mean within-split seed variance.

\subsection{External and survival validation}

External samples were normalized to 10{,}000 counts per spot, log-transformed, restricted to the training gene panel, converted to graphs as above and scored with the seed-7 model (argmax). The MES reference was the mean expression of Neftel MES1/MES2 meta-module markers present in the panel \cite{neftel2019}. For TCGA-GBM, genes were z-scored across samples. Four curated signatures (immune-cold/MES--myeloid, 22 genes; immune-hot/T cell, 18; necrotic/hypoxia, 16; normal brain, 16) were averaged per sample and combined with the pseudo-risk weights. Cox proportional-hazards models (univariate and age-adjusted), the concordance index and Kaplan--Meier curves with log-rank tests at the median were computed with lifelines \cite{davidsonpilon2019}.

\subsection{Software and reproducibility}

The analyses use Python 3.11, PyTorch 2.7 (2.13 for the multi-split and GAT runs), PyTorch Geometric 2.8, scanpy 1.11, scikit-learn, lifelines 0.30 and umap-learn \cite{mcinnes2018}. Numbered scripts reproduce every table and figure from raw data, and unit tests cover graph construction, labeling and models. Figures use the Okabe--Ito color-blind-safe palette.

% ============================================================
\section*{Data and code availability}
Code (MIT license), configuration files, patient splits, trained weights, per-run metrics and all figures and tables are available at \url{https://github.com/hossainlab/deepgraph-gbm}. Training data: SNUH GBM atlas \cite{park2026} (\url{https://cells.ucsc.edu/?ds=multiomic-gbm}). External cohort: Greenwald et al.\ \cite{greenwald2024} (Zenodo 12624108). Survival data: TCGA-GBM via the NCI Genomic Data Commons.

\section*{Competing interests}
The author declares no competing interests.

\section*{Acknowledgments}
We thank the investigators who generated and openly shared the SNUH GBM atlas, the Greenwald et al.\ cohort and the TCGA data.

% ============================================================
\bibliographystyle{unsrtnat}
\bibliography{references}
\clearpage
\appendix
\section*{Supplementary Material}

\renewcommand{\thefigure}{S\arabic{figure}}
\setcounter{figure}{0}
\renewcommand{\thetable}{S\arabic{table}}
\setcounter{table}{0}

\subsection*{Supplementary Figures}

\begin{figure}[H]
  \centering
  \begin{subfigure}[b]{0.32\textwidth}
    \includegraphics[width=\textwidth]{FigureS1a_SNU16B_niche_true.png}
    \caption{True niche}
  \end{subfigure}\hfill
  \begin{subfigure}[b]{0.32\textwidth}
    \includegraphics[width=\textwidth]{FigureS1b_SNU16B_niche_pred.png}
    \caption{Predicted niche}
  \end{subfigure}\hfill
  \begin{subfigure}[b]{0.32\textwidth}
    \includegraphics[width=\textwidth]{FigureS1c_SNU16B_mes.png}
    \caption{MES probability}
  \end{subfigure}
  \caption{\textbf{Spatial niche and MES maps for held-out test section SNU16B.} (a) Silver-standard niche labels; (b) predicted niche labels (seed-7 GraphSAGE, argmax); (c) predicted MES cancer-cell probability per spot.}
  \label{fig:s1}
\end{figure}

\begin{figure}[H]
  \centering
  \begin{subfigure}[b]{0.32\textwidth}
    \includegraphics[width=\textwidth]{FigureS2a_SNU33A_niche_true.png}
    \caption{True niche}
  \end{subfigure}\hfill
  \begin{subfigure}[b]{0.32\textwidth}
    \includegraphics[width=\textwidth]{FigureS2b_SNU33A_niche_pred.png}
    \caption{Predicted niche}
  \end{subfigure}\hfill
  \begin{subfigure}[b]{0.32\textwidth}
    \includegraphics[width=\textwidth]{FigureS2c_SNU33A_mes.png}
    \caption{MES probability}
  \end{subfigure}
  \caption{\textbf{Spatial niche and MES maps for held-out test section SNU33A.} (a) Silver-standard niche labels; (b) predicted niche labels (seed-7 GraphSAGE, argmax); (c) predicted MES cancer-cell probability per spot.}
  \label{fig:s2}
\end{figure}

\begin{figure}[H]
  \centering
  \begin{subfigure}[b]{0.32\textwidth}
    \includegraphics[width=\textwidth]{FigureS3a_SNU33B_niche_true.png}
    \caption{True niche}
  \end{subfigure}\hfill
  \begin{subfigure}[b]{0.32\textwidth}
    \includegraphics[width=\textwidth]{FigureS3b_SNU33B_niche_pred.png}
    \caption{Predicted niche}
  \end{subfigure}\hfill
  \begin{subfigure}[b]{0.32\textwidth}
    \includegraphics[width=\textwidth]{FigureS3c_SNU33B_mes.png}
    \caption{MES probability}
  \end{subfigure}
  \caption{\textbf{Spatial niche and MES maps for held-out test section SNU33B.} (a) Silver-standard niche labels; (b) predicted niche labels (seed-7 GraphSAGE, argmax); (c) predicted MES cancer-cell probability per spot.}
  \label{fig:s3}
\end{figure}

\begin{figure}[H]
  \centering
  \begin{subfigure}[b]{0.32\textwidth}
    \includegraphics[width=\textwidth]{FigureS4a_SNU34A_niche_true.png}
    \caption{True niche}
  \end{subfigure}\hfill
  \begin{subfigure}[b]{0.32\textwidth}
    \includegraphics[width=\textwidth]{FigureS4b_SNU34A_niche_pred.png}
    \caption{Predicted niche}
  \end{subfigure}\hfill
  \begin{subfigure}[b]{0.32\textwidth}
    \includegraphics[width=\textwidth]{FigureS4c_SNU34A_mes.png}
    \caption{MES probability}
  \end{subfigure}
  \caption{\textbf{Spatial niche and MES maps for held-out test section SNU34A.} (a) Silver-standard niche labels; (b) predicted niche labels (seed-7 GraphSAGE, argmax); (c) predicted MES cancer-cell probability per spot.}
  \label{fig:s4}
\end{figure}

\clearpage
\subsection*{Supplementary Tables}

\begin{table}[H]
\centering
\caption{\textbf{Baseline model metrics under default argmax scoring on held-out test patients.} Original, mismatched protocol: single run trained for 120 epochs without early stopping, scored by argmax. Shown for comparison with Table~\ref{tab:comparison}; see Figure~\ref{fig:reversal}.}
\label{tab:s1}
\small
\begin{tabular}{lcc}
\toprule
\textbf{Metric} & \textbf{Multinomial Logistic Regression} & \textbf{Multilayer Perceptron (MLP)} \\
\midrule
Macro-F1          & 0.350 & 0.449 \\
Weighted-F1       & 0.437 & 0.553 \\
Accuracy          & 0.369 & 0.492 \\
AUROC immune-cold & 0.604 & 0.703 \\
AUROC immune-hot  & 0.957 & 0.933 \\
AUROC normal      & 0.572 & 0.698 \\
AUROC necrotic    & 0.782 & 0.854 \\
MES AUROC         & 0.671 & 0.692 \\
MES Pearson $r$   & 0.213 & 0.244 \\
\bottomrule
\end{tabular}
\end{table}

\vspace{1em}

\begin{table}[H]
\centering
\caption{\textbf{External validation per-sample summary on the Greenwald et al.\ cohort} (13 region-labeled Visium samples). Seed-7 GraphSAGE model applied without retraining, argmax decisions; predicted risk from slide hazard head; Pearson correlation $r$ computed against external Neftel MES meta-module expression per spot; predicted niche fractions across spots.}
\label{tab:s2}
\footnotesize
\setlength{\tabcolsep}{4pt}
\begin{tabular}{llccccccc}
\toprule
\textbf{Sample ID} & \textbf{Annotated Region} & \textbf{Spots} & \textbf{Risk} & \textbf{MES $r$} & \textbf{Cold} & \textbf{Hot} & \textbf{Normal} & \textbf{Necrotic} \\
\midrule
GBM\_MGH258      & Unknown            & 2{,}397 & 0.494  & 0.040  & 0.000 & 0.290 & 0.003 & 0.708 \\
GBM\_ZH1007inf   & Infiltrating       & 1{,}436 & 0.418  & $-$0.339 & 0.000 & 0.923 & 0.000 & 0.077 \\
GBM\_ZH1007nec   & Necrotic           & 955   & 0.222  & 0.478  & 0.000 & 0.571 & 0.000 & 0.429 \\
GBM\_ZH1019T1    & Contrast-enhancing & 1{,}112 & 0.283  & 0.096  & 0.000 & 0.266 & 0.000 & 0.734 \\
GBM\_ZH1019inf   & Infiltrating       & 1{,}726 & 0.449  & 0.331  & 0.000 & 0.539 & 0.000 & 0.461 \\
GBM\_ZH8811Abulk & Bulk               & 2{,}486 & 0.025  & 0.577  & 0.000 & 0.362 & 0.000 & 0.638 \\
GBM\_ZH8811Bbulk & Bulk               & 2{,}595 & 0.036  & 0.593  & 0.000 & 0.300 & 0.000 & 0.700 \\
GBM\_ZH8812bulk  & Bulk               & 2{,}343 & 0.575  & 0.623  & 0.093 & 0.684 & 0.004 & 0.219 \\
GBM\_ZH881T1     & Contrast-enhancing & 2{,}792 & 0.002  & 0.322  & 0.000 & 0.340 & 0.000 & 0.660 \\
GBM\_ZH881\_inf  & Infiltrating       & 1{,}387 & 0.433  & 0.162  & 0.000 & 0.133 & 0.076 & 0.790 \\
GBM\_ZH916T1     & Contrast-enhancing & 2{,}780 & 0.078  & 0.298  & 0.000 & 0.588 & 0.000 & 0.412 \\
GBM\_ZH916bulk   & Bulk               & 2{,}199 & $-$0.020 & 0.649 & 0.003 & 0.107 & 0.006 & 0.884 \\
GBM\_ZH916inf    & Infiltrating       & 1{,}280 & 0.316  & 0.012  & 0.265 & 0.148 & 0.174 & 0.412 \\
\bottomrule
\end{tabular}
\end{table}

\clearpage

\begin{table}[H]
\centering
\caption{\textbf{Section-level metadata for the SNUH human glioblastoma atlas} (28 sections from 13 IDH-wildtype patients). Labeled spots are those assigned a niche by the label rules. Slide risk is the rule-derived pseudo-risk target used to train the slide-level head (not a prediction). Split indicates assignment in the primary benchmark split.}
\label{tab:s3}
\footnotesize
\setlength{\tabcolsep}{4.5pt}
\begin{tabular}{llccccc}
\toprule
\textbf{Section} & \textbf{Patient} & \textbf{Tissue Source} & \textbf{Total Spots} & \textbf{Labeled Spots} & \textbf{Slide Risk} & \textbf{Split} \\
\midrule
SNU16A         & SNU16   & FFPE   & 4{,}712 & 4{,}305 & 0.491  & Test \\
SNU16B         & SNU16   & FFPE   & 3{,}853 & 3{,}716 & 0.885  & Test \\
SNU17A         & SNU17   & FFPE   & 2{,}077 & 1{,}965 & 0.652  & Train \\
SNU17B         & SNU17   & FFPE   & 2{,}278 & 2{,}069 & 0.461  & Train \\
SNU18A         & SNU18   & FFPE   & 2{,}633 & 2{,}076 & $-$0.308 & Train \\
SNU18B         & SNU18   & FFPE   & 2{,}638 & 2{,}203 & $-$0.075 & Train \\
SNU18A frozen  & SNU18   & Frozen & 4{,}689 & 2{,}642 & 0.738  & Train \\
SNU21A         & SNU21   & FFPE   & 3{,}722 & 1{,}097 & 0.227  & Train \\
SNU21B         & SNU21   & FFPE   & 4{,}302 & 281     & 0.604  & Train \\
SNU21A frozen  & SNU21   & Frozen & 2{,}408 & 1{,}959 & 0.726  & Train \\
SNU21B frozen  & SNU21   & Frozen & 1{,}572 & 889     & $-$0.772 & Train \\
SNU21C frozen  & SNU21   & Frozen & 3{,}804 & 3{,}558 & $-$0.794 & Train \\
SNU21D frozen  & SNU21   & Frozen & 3{,}467 & 2{,}482 & 0.851  & Train \\
SNU23A         & SNU23   & FFPE   & 4{,}888 & 4{,}061 & $-$0.032 & Train \\
SNU24A         & SNU24   & FFPE   & 3{,}391 & 2{,}921 & 0.989  & Val \\
SNU24B         & SNU24   & FFPE   & 4{,}208 & 3{,}646 & 0.130  & Val \\
SNU25A         & SNU25   & FFPE   & 2{,}610 & 975     & $-$0.344 & Train \\
SNU25B         & SNU25   & FFPE   & 4{,}221 & 3{,}576 & $-$0.412 & Train \\
SNU27A         & SNU27   & FFPE   & 3{,}980 & 3{,}186 & 0.624  & Train \\
SNU27B         & SNU27   & FFPE   & 3{,}851 & 3{,}247 & 0.372  & Train \\
SNU33A         & SNU33   & FFPE   & 3{,}371 & 996     & 0.083  & Test \\
SNU33B         & SNU33   & FFPE   & 3{,}303 & 2{,}591 & 0.990  & Test \\
SNU34A         & SNU34   & FFPE   & 4{,}794 & 4{,}769 & 0.983  & Test \\
SNU43A         & SNU43   & FFPE   & 4{,}701 & 3{,}493 & 0.517  & Train \\
SNU46A         & SNU46   & FFPE   & 2{,}536 & 2{,}452 & $-$0.799 & Train \\
SNU46B         & SNU46   & FFPE   & 3{,}836 & 3{,}713 & $-$0.796 & Train \\
SNU51A         & SNU51   & FFPE   & 4{,}458 & 3{,}539 & 0.922  & Val \\
SNU51B         & SNU51   & FFPE   & 4{,}185 & 3{,}290 & 0.442  & Val \\
\bottomrule
\end{tabular}
\end{table}

\vspace{1em}
\noindent\textbf{Supplementary Table S4.} Spot-level model predictions and external Neftel MES reference scores across all 25{,}488 spots of the Greenwald cohort are provided as a machine-readable dataset in the project repository (\texttt{TableS4\_}\allowbreak\texttt{greenwald\_}\allowbreak\texttt{predictions.csv}, 2.5\,MB).

\vspace{0.5em}
\noindent\textbf{Supplementary Table S5.} Patient-level TCGA niche-composition risk scores and matched clinical outcome covariates are provided as a machine-readable dataset in the project repository (\texttt{TableS5\_}\allowbreak\texttt{tcga\_}\allowbreak\texttt{scores\_}\allowbreak\texttt{clinical.csv}, 26\,KB).

\clearpage

\setcounter{table}{5}
\begin{table}[H]
\centering
\caption{\textbf{Per-seed metrics for all models under the matched protocol on the fixed split.} Seeds 42, 7, 123. All models evaluated with validation-tuned per-class thresholds applied to held-out test patients.}
\label{tab:s6}
\footnotesize
\setlength{\tabcolsep}{3.5pt}
\begin{tabular}{lccccccccccc}
\toprule
\textbf{Model} & \textbf{Seed} & \textbf{Macro-F1} & \textbf{W-F1} & \textbf{Acc} & \textbf{AUC Cold} & \textbf{AUC Hot} & \textbf{AUC Norm} & \textbf{AUC Necr} & \textbf{MES AUC} & \textbf{MES $r$} \\
\midrule
GNN & 42  & 0.458 & 0.641 & 0.705 & 0.661 & 0.975 & 0.693 & 0.943 & 0.794 & 0.375 \\
GNN & 7   & 0.504 & 0.698 & 0.740 & 0.821 & 0.979 & 0.729 & 0.956 & 0.746 & 0.343 \\
GNN & 123 & 0.505 & 0.782 & 0.804 & 0.865 & 0.889 & 0.795 & 0.954 & 0.713 & 0.358 \\
\midrule
MLP & 42  & 0.554 & 0.754 & 0.785 & 0.827 & 0.989 & 0.845 & 0.935 & 0.841 & 0.478 \\
MLP & 7   & 0.593 & 0.771 & 0.793 & 0.851 & 0.986 & 0.831 & 0.948 & 0.809 & 0.436 \\
MLP & 123 & 0.588 & 0.766 & 0.790 & 0.834 & 0.987 & 0.811 & 0.946 & 0.828 & 0.462 \\
\midrule
Logistic & 42  & 0.554 & 0.724 & 0.748 & 0.760 & 0.991 & 0.687 & 0.925 & 0.751 & 0.410 \\
Logistic & 7   & 0.554 & 0.724 & 0.748 & 0.760 & 0.991 & 0.687 & 0.925 & 0.751 & 0.410 \\
Logistic & 123 & 0.554 & 0.724 & 0.748 & 0.760 & 0.991 & 0.687 & 0.925 & 0.751 & 0.410 \\
\bottomrule
\end{tabular}
\end{table}

\vspace{1.5em}

\begin{table}[H]
\centering
\caption{\textbf{Multi-split held-out test macro-F1 across six patient-level splits and three seeds (72 trained models).} Shown as mean $\pm$ standard deviation across seeds 42, 7, and 123 for each split. Logistic regression is deterministic given the split data.}
\label{tab:s7}
\small
\begin{tabular}{lcccc}
\toprule
\textbf{Patient Split} & \textbf{GraphSAGE GNN} & \textbf{GAT (4 heads)} & \textbf{MLP (Matched)} & \textbf{Logistic Regression} \\
\midrule
Split 0 (Original) & 0.489 $\pm$ 0.027 & 0.539 $\pm$ 0.063 & 0.571 $\pm$ 0.021 & 0.554 \\
Split 1            & 0.634 $\pm$ 0.048 & 0.702 $\pm$ 0.039 & 0.690 $\pm$ 0.056 & 0.683 \\
Split 2            & 0.321 $\pm$ 0.032 & 0.355 $\pm$ 0.053 & 0.371 $\pm$ 0.010 & 0.425 \\
Split 3            & 0.412 $\pm$ 0.060 & 0.451 $\pm$ 0.057 & 0.506 $\pm$ 0.014 & 0.530 \\
Split 4            & 0.543 $\pm$ 0.057 & 0.556 $\pm$ 0.091 & 0.556 $\pm$ 0.067 & 0.625 \\
Split 5            & 0.522 $\pm$ 0.094 & 0.538 $\pm$ 0.055 & 0.465 $\pm$ 0.015 & 0.616 \\
\midrule
\textbf{Split Mean $\pm$ SD} & \textbf{0.487 $\pm$ 0.109} & \textbf{0.524 $\pm$ 0.116} & \textbf{0.526 $\pm$ 0.108} & \textbf{0.572 $\pm$ 0.090} \\
\bottomrule
\end{tabular}
\end{table}

\vspace{1.5em}

\begin{table}[H]
\centering
\caption{\textbf{Per-class validation and test metrics for the representative GraphSAGE model} (seed 7, best epoch 12). Per-class thresholds were tuned on validation patients (immune-cold 0.30, immune-hot 0.05, normal 0.725, necrotic 0.15) and applied to test. $n$ is the number of labeled spots in each class. Chance AUPRC equals class prevalence.}
\label{tab:s8}
\small
\begin{tabular}{lcccccc}
\toprule
 & \multicolumn{3}{c}{\textbf{Validation (SNU24, SNU51)}} & \multicolumn{3}{c}{\textbf{Test (SNU16, SNU33, SNU34)}} \\
\cmidrule(lr){2-4}\cmidrule(lr){5-7}
\textbf{Class} & $n$ & \textbf{AUROC / AUPRC} & \textbf{F1} & $n$ & \textbf{AUROC / AUPRC} & \textbf{F1} \\
\midrule
Immune-cold & 8{,}638 & 0.866 / 0.904 & 0.874 & 11{,}125 & 0.821 / 0.913 & 0.836 \\
Immune-hot  & 15      & 0.786 / 0.003 & 0.012 & 185     & 0.979 / 0.573 & 0.473 \\
Normal      & 2{,}715 & 0.977 / 0.917 & 0.847 & 1{,}455 & 0.729 / 0.266 & 0.241 \\
Necrotic    & 2{,}028 & 0.902 / 0.738 & 0.604 & 3{,}612 & 0.956 / 0.862 & 0.466 \\
\midrule
Macro-F1 / accuracy & \multicolumn{3}{c}{0.584 / 0.827} & \multicolumn{3}{c}{0.504 / 0.740} \\
MES AUROC / AUPRC / $r$ & \multicolumn{3}{c}{0.795 / 0.438 / 0.400} & \multicolumn{3}{c}{0.746 / 0.388 / 0.343} \\
\bottomrule
\end{tabular}
\end{table}

\vspace{1.5em}

\begin{table}[H]
\centering
\caption{\textbf{Paired comparisons across six patient-level splits.} Differences in split-mean macro-F1 (model A $-$ model B; each split mean averages three seeds). $t$-test df $= 5$. The minimum attainable two-sided Wilcoxon $p$ with $n = 6$ is 0.031.}
\label{tab:s9}
\small
\begin{tabular}{lccccc}
\toprule
\textbf{A vs.\ B} & \textbf{Mean diff.\ (95\% CI)} & $t$ & $p_t$ & $p_\mathrm{Wilcoxon}$ & \textbf{A wins} \\
\midrule
Logistic vs.\ GraphSAGE & 0.086 (0.059, 0.112)   & 8.27  & 0.0004 & 0.031 & 6/6 \\
MLP vs.\ GraphSAGE      & 0.040 ($-$0.018, 0.098) & 1.77  & 0.137  & 0.219 & 5/6 \\
GAT vs.\ GraphSAGE      & 0.037 (0.015, 0.059)   & 4.36  & 0.007  & 0.031 & 6/6 \\
GAT vs.\ MLP            & $-$0.003 ($-$0.049, 0.044) & $-$0.15 & 0.886 & 0.844 & 3/6 \\
GAT vs.\ Logistic       & $-$0.049 ($-$0.092, $-$0.005) & $-$2.90 & 0.034 & 0.094 & 1/6 \\
Logistic vs.\ MLP       & 0.046 ($-$0.019, 0.110) & 1.82  & 0.128  & 0.156 & 4/6 \\
\bottomrule
\end{tabular}
\end{table}

\end{document}
